\documentclass[11pt]{article}
\usepackage{graphicx}
\usepackage{amsmath}
\usepackage{booktabs}
\usepackage[margin=1in]{geometry}
\usepackage{setspace}
\usepackage{xcolor}
\usepackage{hyperref}
\hypersetup{colorlinks=true, linkcolor=blue, urlcolor=blue}

\title{License-Cap Policy Evaluation: A Synthesis\\
\large Incidence, Magnitude, and the Politics of the Local-Option Margin}
\author{Working notes -- entry game-marijuana project}
\date{August 2026}

\begin{document}
\maketitle

\begin{abstract}
\noindent This note assembles a policy-evaluation section from results already estimated across
Colorado, Washington, and Montana, plus one new test. Three points are established. First,
\textbf{incidence}: where cannabis license entry is restricted (Colorado's local-option bans,
Washington's statewide numerical cap), a cross-border demand shock lands on revenue per existing
dispensary rather than on the number of dispensaries or revenue per capita; this holds not just in the
pooled Colorado-Washington estimate but in each state's licensing regime tested independently. Where
entry is not capped by a fixed number (Montana, under a moratorium restricted to already-licensed 2020
medical operators), the same shock shows up on revenue per capita instead. Second, \textbf{magnitude}:
using each state's best-supported coefficient, the capped regime captures roughly 88\% (\$164M) of the
combined 2025 border-effect premium and Montana 12\% (\$23M); normalized by the population of nearby
non-legal-state residents each regime could plausibly draw from, the capped regime still extracts about
2.7 times more per potential customer, so this is not simply a market-size artifact. Third, a genuinely
new result: Colorado's border counties are, if anything, \emph{less} likely to ever have permitted an
operating dispensary than interior counties (43\% vs.\ 72\%, $p=0.036$), the opposite of what a simple
fiscal-capture theory of local opt-in decisions would predict. This is reported as an open, unresolved
finding, not smoothed into the main story, and is discussed as a direction for a genuinely new
contribution rather than a robustness check.
\end{abstract}

\tableofcontents

\section{Framing}

The paper's core empirical pattern is that dispensaries in counties near a non-legal-state border earn
more revenue than interior dispensaries of the same state, and that this shows up on different margins
depending on how restrictive each state's entry policy is. This note turns that pattern into an explicit
policy-evaluation section, organized around three questions a referee is likely to ask:

\begin{enumerate}
\item \textbf{Does the incidence result generalize across different licensing mechanisms}, or is it an
artifact of one particular policy design (e.g., Washington's specific numerical cap)?
\item \textbf{How large is the transfer in dollar terms}, and is the difference across regimes just a
market-size effect?
\item \textbf{Is there a political-economy mechanism behind the local-option version of the cap}
(Colorado), where local governments choose whether to allow retail marijuana at all?
\end{enumerate}

Consistent with the rest of the project, none of what follows is presented as a structural welfare
estimate. The revenue data contain no prices or quantities, so nothing here can separate a genuine
increase in served demand from market power at capped, scarce locations. The contribution is a clean,
falsifiable description of \emph{where} a real, sizable demand shock lands under different licensing
regimes, and a preliminary look at \emph{why} local governments themselves might restrict entry.

\section{Part 1: Incidence Generalizes Across Cap Mechanisms}
\label{sec:incidence}

Colorado and Washington restrict entry through different institutional channels. Washington caps the
statewide number of retail licenses (334 initially, raised to 556 in 2016). Colorado imposes no statewide
number at all; instead, Amendment 64 gives every county and municipality the authority to ban retail
marijuana outright, and most have. If the incidence result, effect concentrated in revenue per
dispensary, not in entry or revenue per capita, is really about entry restriction as a general
phenomenon, it should hold independently in \emph{both} states, not just in a pooled estimate that could
be driven by one of them.

\begin{table}[htbp]
\centering
\caption{Incidence of the Border Effect, by State and Margin}
\label{tab:incidence}
\small
\begin{tabular}{lccc}
\toprule
& Colorado & Washington & Montana \\
& (local-option bans) & (statewide cap, 334$\to$556) & (moratorium, 2020 medical only) \\
\midrule
Entry density (disp./1{,}000) & $+0.015$ & $-0.002$ & $+0.142^{a}$ \\
& ($p=0.889$) & ($p=0.940$) & ($p=0.106$) \\
Revenue per dispensary & $+0.552^{**}$ & $+0.233$ & $+0.012$ \\
& ($p=0.005$) & ($p=0.334$) & ($p=0.958$) \\
Revenue per capita & $-0.013$ & $+0.341^{b}$ & $+0.654^{**}$ \\
& ($p=0.968$) & ($p=0.121$) & ($p=0.036$) \\
\bottomrule
\end{tabular}
\begin{minipage}{0.95\textwidth}
\vspace{0.3em}
\footnotesize
\textit{Notes:} All coefficients from the same specification used throughout the paper (log outcome,
border indicator, log population, year fixed effects where a panel), estimated separately by state
rather than pooled. Entry density and revenue per dispensary use 2025 cross-sections; revenue per capita
uses each state's full panel. $p$-values are county-clustered where a panel exists, HC1 for single-year
cross-sections. $^{a}$Montana's entry-density coefficient is significant without the population control
($p=0.035$) and marginal with it. $^{b}$Washington's revenue-per-capita coefficient is significant under
HC1 ($p=0.001$) but not under county clustering; reported as fragile, not null. $^{**}p<0.05$.
\end{minipage}
\end{table}


\begin{figure}[htbp]
\centering
\includegraphics[width=0.95\textwidth]{synthesis_fig_incidence.png}
\caption{Border coefficient on three outcomes, each state's licensing regime tested independently.
Entry density and revenue per dispensary are estimated the same way throughout the paper; revenue per
capita uses the same specification, run separately by state rather than pooled.}
\label{fig:incidence}
\end{figure}

Table~\ref{tab:incidence} and Figure~\ref{fig:incidence} show the result holds up. In Colorado, entry
density is a precise null ($b=+0.015$, $p=0.89$) and revenue per dispensary is large and significant
($b=+0.552$, $p=0.005$). In Washington, entry density is likewise a precise null ($b=-0.002$, $p=0.94$)
and revenue per dispensary is positive, though not significant on its own ($b=+0.233$, $p=0.33$). Both
states independently show the same qualitative pattern that the pooled estimate shows: no detectable
entry response, and a demand shock that concentrates in revenue per existing store. This is the single
most useful fact for a referee, the result is not an artifact of Washington's specific 334/556-license
cap; a completely different institutional mechanism, Colorado's local-option ban authority, produces the
same signature.

One nuance worth stating directly: Washington's revenue-per-capita coefficient ($b=+0.341$) is
significant under HC1 ($p=0.001$) but not under county clustering ($p=0.12$), so Washington shows a
weak, statistically fragile hint of a per-capita response that Colorado does not share at all
($b=-0.013$, $p=0.97$). This is consistent with Washington's licensing regime being formally a numerical
cap but administered with periodic capacity increases (334 to 556 in 2016), giving it slightly more
flexibility than Colorado's local-option bans, which, once a jurisdiction opts out, do not expire on
their own. Montana, with no numerical cap of either kind, shows the reverse signature throughout: entry
density positive and marginal ($b=+0.142$, $p=0.11$; significant without controls, $p=0.035$), revenue
per dispensary a precise null ($b=+0.012$, $p=0.96$), and revenue per capita significant
($b=+0.654$, $p=0.036$).

\section{Part 2: Magnitude, and Whether It Is Just Market Size}
\label{sec:magnitude}

\begin{table}[htbp]
\centering
\caption{Fiscal Transfer, Capped vs.\ Moratorium Regime, 2025}
\label{tab:magnitude}
\small
\begin{tabular}{lcc}
\toprule
& Colorado + Washington & Montana \\
& (capped regime) & (moratorium regime) \\
\midrule
Coefficient used & $+0.614$ & $+0.654$ \\
& (2025-specific, $p=0.002$) & (pooled 2022--2025, $p=0.036$) \\
2025 border-county revenue & \$357.5M & \$48.3M \\
\textbf{Attributable to border effect} & \textbf{\$164.0M} & \textbf{\$23.2M} \\
Share of combined premium (\$187.2M) & 87.6\% & 12.4\% \\
Share of border-county revenue & 45.9\% & 48.0\% \\
Catchment population$^{a}$ & 512{,}162 & 198{,}016 \\
\textbf{Transferred per catchment resident} & \textbf{\$320.16} & \textbf{\$117.19} \\
\bottomrule
\end{tabular}
\begin{minipage}{0.95\textwidth}
\vspace{0.3em}
\footnotesize
\textit{Notes:} $^{a}$Catchment population is the union of non-legal-state counties within 70 miles of any
border county in that regime (20 counties across UT/WY/NE/ID for Colorado+Washington; 10 counties across
ID/ND/WY for Montana), not a per-county sum, to avoid double-counting non-legal counties near more than
one border county. Attributable revenue $= \text{actual border revenue} \times \text{premium}/(1+\text{premium})$,
where premium $=100(e^{\beta}-1)$. This is an accounting decomposition of an already-estimated
coefficient, not an independent causal magnitude.
\end{minipage}
\end{table}


\begin{figure}[htbp]
\centering
\includegraphics[width=0.95\textwidth]{visualized_premium_share.png}
\caption{Left: share of the combined 2025 border-effect premium captured by each regime. Right: the same
premium normalized by catchment population, the union of non-legal-state counties within 70 miles of any
of that regime's border counties, not a per-county sum, to avoid double-counting overlapping catchments.}
\label{fig:magnitude}
\end{figure}

Table~\ref{tab:magnitude} and Figure~\ref{fig:magnitude} report the dollar magnitude, using each regime's
own best-supported coefficient: Colorado and Washington's 2025-specific coefficient ($b=0.614$,
$p=0.002$, well-powered given a full decade of data and 12 border counties), and Montana's pooled
2022--2025 coefficient ($b=0.654$, $p=0.036$), since Montana's 2025-only estimate is not significant on
its own (only 8 border counties in one year). Of the \$187.2M combined premium, the capped regime
captures 88\% (\$164M) and Montana 12\% (\$23M).

Raw dollars conflate two things: a genuinely different regime effect, and the simple fact that Colorado
and Washington's cannabis economy is much larger than Montana's. To separate them, we computed each
regime's \emph{catchment population}, the union of non-legal-state counties within 70 miles of any of
that regime's own border counties (20 counties across UT/WY/NE/ID for Colorado+Washington, 512{,}162
residents; 10 counties across ID/ND/WY for Montana, 198{,}016 residents). Catchment population is 2.6
times larger for the capped regime, but the dollar transfer is 4.45 times larger, so market size explains
part, but not all, of the gap. Normalized per catchment resident, the capped regime still extracts about
\$320 versus Montana's \$117, roughly 2.7 times more per potential cross-border customer. The capped
regime is not simply bigger, it appears to be more effective at converting a given pool of nearby
non-legal-state residents into revenue, consistent with the idea that restricting entry concentrates
demand onto fewer, more capitalized locations rather than diffusing it across many marginal ones.

We do not read this as a pure regime effect. Colorado and Washington's markets are also roughly eight
years older than Montana's (2014/2015 launch vs.\ 2022), and the paper's own market-age event study found
the Colorado-Washington premium was not reliably significant until roughly year 7 of a state's market
life. Montana's own within-state trend across its four available years does not show the same
build-up, its per-capita border gap is roughly flat (58\% in 2022, 65\% in 2023, 55\% in 2024, 64\% in
2025), so market immaturity alone does not cleanly explain the gap either. Untangling market age from
regime type would need more years of Montana data than currently exist.

\section{Part 3: A Falsifiable Prediction, Not Just a Pattern}
\label{sec:prediction}

What separates a policy-evaluation section from a restatement of results is a claim the data could have
contradicted. The prediction here is explicit: \emph{under a binding entry restriction, a cross-border
demand shock cannot expand the number of sellers, so it must appear as higher revenue per existing
seller; under unrestricted entry, the same shock should instead expand the number of sellers, and revenue
per seller should not need to rise as much.} This is not fit to the data after the fact, it is the same
logic already in the paper's conceptual framework (Section~3, citing Bresnahan and Reiss's classic
entry-threshold result), and Montana was chosen as a test case specifically because it is the one sample
state without a statewide numerical cap, before its entry-density and revenue-per-capita results were
run.

The prediction could have failed in an observable way, Montana could have shown the same
per-dispensary-only pattern as Colorado and Washington, which would have undercut the license-cap
interpretation entirely. It did not: Montana's reversal (per-capita significant, per-dispensary null) is
exactly the mirror image the prediction implies. The New Mexico natural experiment (paper
Section~6.1) adds a second, independent falsification opportunity on a different margin: if the premium
is capacity-constrained demand rather than a stable, unobserved difference between border and interior
counties, it should fall when a neighboring county's non-legal-state neighbor itself legalizes. It does,
by 38.6\% in Colorado counties newly bordering New Mexico, precisely timed to New Mexico's own 2022
retail opening, not before.

\section{Part 4: A New Result, and an Honest Reversal}
\label{sec:politicaleconomy}

Colorado's cap is not really a cap at all, it is a collection of independent local opt-in/opt-out
decisions under Amendment 64. This opens a question the paper has not tested: \textbf{do local
governments in border counties, who stand to capture the fiscal windfall documented in
Sections~\ref{sec:incidence}--\ref{sec:magnitude}, allow retail marijuana at a higher rate than interior
counties?} A pure fiscal-competition story predicts yes.

\begin{table}[htbp]
\centering
\caption{Do Colorado Border Counties Opt In to Retail Marijuana at a Higher Rate?}
\label{tab:politicaleconomy}
\small
\begin{tabular}{lccc}
\toprule
& Interior counties & Border counties & \\
& ($n=50$) & ($n=14$) & \\
\midrule
Share ever with $\geq$1 dispensary & 72.0\% & 42.9\% & \\
\midrule
Specification & Coefficient & $p$-value & $n$ \\
\midrule
No controls & $-0.291$ & $0.051$ & 64 \\
+ log(population) & $-0.263^{**}$ & $0.036$ & 64 \\
+ log(population), pop.\ $\geq$5{,}000 only & $-0.258$ & $0.105$ & 52 \\
\bottomrule
\end{tabular}
\begin{minipage}{0.95\textwidth}
\vspace{0.3em}
\footnotesize
\textit{Notes:} Outcome is a binary indicator for whether a Colorado county ever had an operating
dispensary between 2015 and 2025, the best available proxy for effective local exclusion of retail
marijuana under Amendment~64's opt-out authority (no compiled municipal-ordinance dataset exists; this
proxy conflates active local bans with no operator ever entering a low-demand market). Sample is the full
64-county Colorado universe, not only the 42 counties in the paper's revenue panel. Linear probability
model, HC1 standard errors. \textbf{The sign is opposite to the fiscal-capture prediction}: border
counties are less, not more, likely to have ever permitted retail marijuana. $^{**}p<0.05$.
\end{minipage}
\end{table}


\begin{figure}[htbp]
\centering
\includegraphics[width=0.68\textwidth]{synthesis_fig_politicaleconomy.png}
\caption{Share of Colorado's 64 counties that ever had an operating dispensary between 2015 and 2025,
border vs.\ interior.}
\label{fig:politicaleconomy}
\end{figure}

\textbf{Data and limitation, stated up front.} No compiled, county-level dataset of actual municipal
marijuana-ban ordinances exists for Colorado; opt-out decisions are made city by city, not uniformly by
county, and building one from scratch would require manually compiling ordinance records for dozens of
jurisdictions, the kind of labor-intensive reconstruction this project has elsewhere judged not worth
attempting without a clearer payoff. The proxy used here instead is whether a county ever had \emph{any}
operating dispensary across the full 2015--2025 panel. This is imperfect: a county with zero dispensaries
could reflect an active local ban, or simply no operator choosing to enter a small, low-demand market even
where retail is legal. It is directionally informative, not a direct measure of ordinance status, and
should be described that way in any writeup.

Table~\ref{tab:politicaleconomy} and Figure~\ref{fig:politicaleconomy} report the result, using the full
64-county Colorado universe (not just the 42 counties already in the revenue panel, which by construction
excludes the counties this test needs). \textbf{The finding is the opposite of the fiscal-capture
prediction.} Only 43\% of Colorado's 14 border counties (within 70 miles of Utah, Wyoming, or Nebraska)
ever had an operating dispensary, versus 72\% of the 50 interior counties, a gap that is significant with
a population control ($b=-0.263$, $p=0.036$) and marginal without one ($p=0.051$). Restricting to counties
with at least 5{,}000 residents, to rule out the smallest, commercially thinnest rural counties, shrinks
the gap somewhat (56\% vs.\ 79\%) and weakens its significance ($p=0.105$), but the direction does not
reverse.

\textbf{We report this as a genuine, unresolved finding, not something to explain away.} A plausible
account is that Colorado's non-legal-state neighbors, Utah, Wyoming, and Nebraska, border its Eastern
Plains and southwestern counties, a more rural, culturally conservative part of the state where local
opposition to retail marijuana plausibly runs higher regardless of the fiscal upside, meaning local
political preference dominates fiscal self-interest rather than the reverse. That account is speculative
and not tested here; it would need actual voting or demographic data on these specific counties, not just
the border-status regression. What can be said with the current test is narrower but still useful: the
data reject the simple version of the fiscal-capture story, and whatever is determining Colorado
municipalities' opt-in decisions, it is not primarily "does this county sit next to a captive
non-legal-market population."

\section{Discussion: Assembling This Into the Paper}

Sections~\ref{sec:incidence}--\ref{sec:prediction} require no new empirical work; they are a synthesis of
results already estimated, organized around the three questions a referee is most likely to raise, and
strengthened by two calculations that were not previously in the paper (Colorado- and Washington-specific
entry density, and Colorado- and Washington-specific revenue per capita). These fill a real gap: the
paper previously only showed the entry-density null pooled, which leaves open the objection that it is
specific to Washington's cap design. That objection is now closed.

Section~\ref{sec:politicaleconomy} is different in kind. It is a genuinely new empirical result, not a
repackaging, and it does not confirm the story one might have hoped for. That is, on balance, a point in
its favor for publication: reviewers are more persuaded by a paper that reports a result it did not
expect and explains honestly what it can and cannot conclude than by one that only reports confirming
evidence. Two honest paths forward with it:

\begin{enumerate}
\item \textbf{Report it as a limitation-flagged secondary finding}, exactly as written above, a real
pattern in the data, with the has-any-dispensary proxy's limitations stated plainly, and no causal claim
about local political preferences beyond noting it as the most plausible untested explanation.
\item \textbf{Invest in actually collecting municipal ordinance or vote-share data} for Colorado's border
region specifically, which would let the paper distinguish "local political opposition" from "no operator
ever tried" and would turn this from a suggestive reversal into a real political-economy result. This is
real additional work, not something to promise without doing it.
\end{enumerate}

\section{Caveats Carried Through From the Rest of the Project}

\begin{itemize}
\item All dollar figures here are accounting decompositions of an already-estimated coefficient, not
independent causal magnitudes; they cannot distinguish genuine demand absorption from market power at
capped locations, consistent with the paper's existing Discussion section and Hollenbeck and Uetake's
(2021) market-power finding for Washington specifically.
\item The core revenue analysis still rests on two states with usable administrative data (Colorado,
Washington) plus Montana as a single institutional contrast; none of the calculations here add a fourth
state.
\item Montana's revenue-per-dispensary test remains a single 2025 cross-section (no historical dispensary
open-dates exist in Montana's registry), materially weaker than its revenue-per-capita panel; this
asymmetry is inherited by every Montana number in this note.
\item The political-economy test's has-any-dispensary proxy conflates active local bans with no operator
ever entering; treat Section~\ref{sec:politicaleconomy} as suggestive, not as evidence of specific
municipal ordinance behavior.
\end{itemize}

\end{document}
